\section{Intertwiner with the space of symmetric traceless tensors}

Let \(\mathcal V\subset S^2\) be the set of the twelve unit vertices
of the icosahedron. For each \(v\in\mathcal V\), define
\begin{equation}
\label{rm:eq:gravity-Qv}
Q_v=vv^T-\frac13I_3\in
\operatorname{STF}_3:=\operatorname{Sym}^2_0(\R^3).
\end{equation}
The frame map
\begin{equation}
\label{rm:eq:gravity-Gamma0}
\Gamma_0\!\left(\sum_{v\in\mathcal V}x_ve_v\right)
=\sum_{v\in\mathcal V}x_vQ_v
\end{equation}
is \(A_5\)-equivariant and satisfies
\begin{equation}
\label{rm:eq:gravity-tight-frame}
\Gamma_0\Gamma_0^*=\frac85I_{\operatorname{STF}_3},
\qquad
\Gamma_0=\Gamma_0P_5.
\end{equation}

\begin{theorem}[Normalized tensor entangler]
\label{rm:thm:gravity-gamma5}
The map
\begin{equation}
\label{rm:eq:gravity-gamma5}
\boxed{
\Gamma_5:=\sqrt{\frac58}\,\Gamma_0|_{V_5}:
V_5\xrightarrow{\ \simeq\ }\operatorname{STF}_3}
\end{equation}
is a \(A_5\)-equivariant isometry. Its adjoint is
\begin{equation}
\label{rm:eq:gravity-gamma5-adjoint}
\Gamma_5^*Q=\sqrt{\frac58}
\sum_{v\in\mathcal V}(v^TQv)e_v.
\end{equation}
\end{theorem}

\begin{proof}
Equivariance follows from \(Q_{gv}=gQ_vg^T\). The operator
\(\Gamma_0\Gamma_0^*\) commutes with the irreducible representation of
\(A_5\) on \(\operatorname{STF}_3\), and is therefore scalar. Its trace
is
\[
\sum_{v\in\mathcal V}\|Q_v\|_F^2
=12\cdot\frac23=8;
\]
dividing by dimension five yields \(8/5\). Hence,
\(\Gamma_5\Gamma_5^*=I\). Since domain and codomain have dimension
five, also \(\Gamma_5^*\Gamma_5=I\). The formula for the adjoint
follows from \(\langle Q,Q_v\rangle_F=v^TQv\).
\end{proof}

Normalization through the icosahedral frame
eliminates a purely cardinal coincidence: it constructs the map that
transports the finite summand to the continuous spin-two
carrier.

\section{Transverse traceless projection}

For \(n\in S^2\), let \(\Pi(n)=I-nn^T\). We define
\begin{equation}
\label{rm:eq:gravity-lambda-tt}
\boxed{
\Lambda(n)Q
=\Pi Q\Pi-\frac12\Tr(\Pi Q\Pi)\Pi.}
\end{equation}

\begin{theorem}[Reduction from five components to two]
\label{rm:thm:gravity-tt}
The restriction of \(\Lambda(n)\) to \(\operatorname{STF}_3\) is the
orthogonal projector onto
\begin{equation}
\label{rm:eq:gravity-htt}
\mathcal H_{\rm TT}(n)
=\{Q\in\operatorname{STF}_3:Qn=0\}.
\end{equation}
In particular,
\begin{equation}
\label{rm:eq:gravity-lambda-properties}
\Lambda(n)^2=\Lambda(n)=\Lambda(n)^*,
\qquad \rank\Lambda(n)=2.
\end{equation}
For \(n=e_3\), the image is generated by
\begin{equation}
\label{rm:eq:gravity-plus-cross}
E_+=\frac1{\sqrt2}\operatorname{diag}(1,-1,0),
\qquad
E_\times=\frac1{\sqrt2}
\begin{pmatrix}0&1&0\\1&0&0\\0&0&0\end{pmatrix}.
\end{equation}
\end{theorem}

\begin{proof}
Since \(\Pi n=0\), the output is transverse. Its trace is
\[
\Tr(\Pi Q\Pi)-\frac12\Tr(\Pi)\Tr(\Pi Q\Pi)=0,
\]
since \(\Tr\Pi=2\). If \(Qn=0\) and \(\Tr Q=0\), then
\(\Pi Q\Pi=Q\) and \(\Tr(\Pi Q\Pi)=0\), so that
\(\Lambda(n)Q=Q\). The formula is self-adjoint with respect to the
Frobenius product. In the frame \(n=e_3\), its matrix in an adapted STF basis is
\(\operatorname{diag}(1,1,0,0,0)\), which proves the rank.
\end{proof}

The complete chain thus remains,
\begin{equation}
\label{rm:eq:gravity-12-5-5-2}
\boxed{
V_{12}\xrightarrow{P_5}V_5
\xrightarrow{\Gamma_5}\operatorname{STF}_3
\xrightarrow{\Lambda(n)}\mathcal H_{\rm TT}(n),
\qquad 12\longrightarrow5\longrightarrow5\longrightarrow2.}
\end{equation}
The reduced coupling
\begin{equation}
\label{rm:eq:gravity-effective-map}
\mathbb G_{{\rm TT},a}(n)
=G_a(\theta_{108}^{\rm circ})\Lambda(n)\Gamma_5P_5
\end{equation}
has rank two.
